% GENERATED FILE. Edit canonical lessons, then run npm run generate:books.
% canonical-source-hash: fnv1a64:7af6fedfcdf2f01c
% canonical-lessons: KA-C294-kadalekayi, KA-C294-masale, KA-C294-catni, KA-C294-duddu, KA-C294-cillare

\chapter{Peanut, Spices, Chutney, Money, Small change}
\label{ch:ka-peanut-spices-chutney-money-small-change}
\hlchaptermodality{294}

\begin{chapteropening}
\textbf{By the end of this chapter:} \emph{I can say a peanut, spices, chutney, money and small change.}

Complete the last lesson of chapter 294: I can say a peanut, spices, chutney, money and small change.
\end{chapteropening}

\section[\texorpdfstring{\kn{ಕಡಲೆಕಾಯಿ} (kaḍalekāyi)}{ಕಡಲೆಕಾಯಿ (kaḍalekāyi)}]{\kn{ಕಡಲೆಕಾಯಿ} (kaḍalekāyi) — a peanut}

\label{lesson:KA-C294-kadalekayi}



\begin{quote}
Before the new one: say the Kannada for a tomato, then the Kannada for a carrot.
\end{quote}

\subsection*{You'll want to know: \kn{ಕಡಲೆಕಾಯಿ}}
\textbf{\kn{ಕಡಲೆಕಾಯಿ}} — \emph{kaḍalekāyi} — \textquotedblleft{}a peanut\textquotedblright{}.

\begin{grammarlens}[title={What you've built}]
One more word for this chapter.
\end{grammarlens}

\subsection*{Your turn}
\begin{itemize}
\raggedright
  \item \emph{Say it:} \emph{kaḍalekāyi}
  \item \emph{Say it:} \emph{kaḍalekāyi}, once more
  \item \emph{Say it:} say \emph{kaḍalekāyi}
  \item \emph{Recall:} say the Kannada for a tomato, then the Kannada for a carrot, then say \emph{kaḍalekāyi} again
\end{itemize}

\begin{tcolorbox}[breakable,colback=teal!4,colframe=teal!35!black,title={Before you move on}]
What does \kn{ಕಡಲೆಕಾಯಿ} mean? (\textquotedblleft{}A peanut\textquotedblright{}.) Say it once more.
\end{tcolorbox}
\section[\texorpdfstring{\kn{ಮಸಾಲೆ} (masāle)}{ಮಸಾಲೆ (masāle)}]{\kn{ಮಸಾಲೆ} (masāle) — spices, a spice mix}

\label{lesson:KA-C294-masale}



\begin{quote}
Before the new one: say the Kannada for a carrot, then the Kannada for a peanut.
\end{quote}

\subsection*{You'll want to know: \kn{ಮಸಾಲೆ}}
\textbf{\kn{ಮಸಾಲೆ}} — \emph{masāle} — \textquotedblleft{}spices, a spice mix\textquotedblright{}.

\begin{grammarlens}[title={What you've built}]
One more word for this chapter.
\end{grammarlens}

\subsection*{Your turn}
\begin{itemize}
\raggedright
  \item \emph{Say it:} \emph{masāle}
  \item \emph{Say it:} \emph{masāle}, once more
  \item \emph{Say it:} say \emph{masāle}
  \item \emph{Recall:} say the Kannada for a carrot, then the Kannada for a peanut, then say \emph{masāle} again
\end{itemize}

\begin{tcolorbox}[breakable,colback=teal!4,colframe=teal!35!black,title={Before you move on}]
What does \kn{ಮಸಾಲೆ} mean? (\textquotedblleft{}Spices, a spice mix\textquotedblright{}.) Say it once more.
\end{tcolorbox}
\section[\texorpdfstring{\kn{ಚಟ್ನಿ} (caṭni)}{ಚಟ್ನಿ (caṭni)}]{\kn{ಚಟ್ನಿ} (caṭni) — chutney}

\label{lesson:KA-C294-catni}



\begin{quote}
Before the new one: say the Kannada for a peanut, then the Kannada for spices.
\end{quote}

\subsection*{You'll want to know: \kn{ಚಟ್ನಿ}}
\textbf{\kn{ಚಟ್ನಿ}} — \emph{caṭni} — \textquotedblleft{}chutney\textquotedblright{}.

\begin{grammarlens}[title={What you've built}]
One more word for this chapter.
\end{grammarlens}

\subsection*{Your turn}
\begin{itemize}
\raggedright
  \item \emph{Say it:} \emph{caṭni}
  \item \emph{Say it:} \emph{caṭni}, once more
  \item \emph{Say it:} say \emph{caṭni}
  \item \emph{Recall:} say the Kannada for a peanut, then the Kannada for spices, then say \emph{caṭni} again
\end{itemize}

\begin{tcolorbox}[breakable,colback=teal!4,colframe=teal!35!black,title={Before you move on}]
What does \kn{ಚಟ್ನಿ} mean? (\textquotedblleft{}Chutney\textquotedblright{}.) Say it once more.
\end{tcolorbox}
\section[\texorpdfstring{\kn{ದುಡ್ಡು} (duḍḍu)}{ದುಡ್ಡು (duḍḍu)}]{\kn{ದುಡ್ಡು} (duḍḍu) — money}

\label{lesson:KA-C294-duddu}



\begin{quote}
Before the new one: say the Kannada for spices, then the Kannada for chutney.
\end{quote}

\subsection*{You'll want to know: \kn{ದುಡ್ಡು}}
\textbf{\kn{ದುಡ್ಡು}} — \emph{duḍḍu} — \textquotedblleft{}money\textquotedblright{}.

\begin{grammarlens}[title={What you've built}]
One more word for this chapter.
\end{grammarlens}

\subsection*{Your turn}
\begin{itemize}
\raggedright
  \item \emph{Say it:} \emph{duḍḍu}
  \item \emph{Say it:} \emph{duḍḍu}, once more
  \item \emph{Say it:} say \emph{duḍḍu}
  \item \emph{Recall:} say the Kannada for spices, then the Kannada for chutney, then say \emph{duḍḍu} again
\end{itemize}

\begin{tcolorbox}[breakable,colback=teal!4,colframe=teal!35!black,title={Before you move on}]
What does \kn{ದುಡ್ಡು} mean? (\textquotedblleft{}Money\textquotedblright{}.) Say it once more.
\end{tcolorbox}
\section[\texorpdfstring{\kn{ಚಿಲ್ಲರೆ} (cillare)}{ಚಿಲ್ಲರೆ (cillare)}]{\kn{ಚಿಲ್ಲರೆ} (cillare) — small change}

\label{lesson:KA-C294-cillare}



\begin{quote}
Before the new one: say the Kannada for chutney, then the Kannada for money.
\end{quote}

\subsection*{You'll want to know: \kn{ಚಿಲ್ಲರೆ}}
\textbf{\kn{ಚಿಲ್ಲರೆ}} — \emph{cillare} — \textquotedblleft{}small change\textquotedblright{}.

\begin{grammarlens}[title={What you've built}]
One more word for this chapter.
\end{grammarlens}

\subsection*{Your turn}
\begin{itemize}
\raggedright
  \item \emph{Say it:} \emph{cillare}
  \item \emph{Say it:} \emph{cillare}, once more
  \item \emph{Say it:} say \emph{cillare}
  \item \emph{Recall:} say the Kannada for chutney, then the Kannada for money, then say \emph{cillare} again
\end{itemize}

\begin{tcolorbox}[breakable,colback=teal!4,colframe=teal!35!black,title={Before you move on}]
What does \kn{ಚಿಲ್ಲರೆ} mean? (\textquotedblleft{}Small change\textquotedblright{}.) Say it once more.
\end{tcolorbox}
